\documentclass[11pt]{article}

\usepackage[margin=1in]{geometry}
\usepackage{amsmath, amssymb}
\usepackage{xcolor}
\usepackage{listings}
\usepackage{hyperref}
\usepackage{array}
\usepackage{booktabs}

\hypersetup{
    colorlinks=true,
    linkcolor=blue!60!black,
    urlcolor=blue!60!black,
    citecolor=blue!60!black
}

\lstdefinestyle{bugpython}{
  language=Python,
  basicstyle=\ttfamily\small,
  keywordstyle=\color{blue!60!black},
  stringstyle=\color{green!40!black},
  commentstyle=\color{gray!70!black},
  showstringspaces=false,
  breaklines=true,
  frame=single,
  rulecolor=\color{gray!30},
  columns=fullflexible,
  keepspaces=true
}

\title{SymPy Bug Report}
\author{}
\date{}

\begin{document}

\maketitle

\section{Bug 51: False solution for atan2(x, 1) = pi in solveset}

\subsection{Minimal Reproducer}

\medskip

\begin{lstlisting}[style=bugpython]
from sympy import Eq, S, atan2, pi, solveset, symbols

x = symbols("x")
eq = Eq(atan2(x, 1), pi)
sol = solveset(eq, x, S.Reals)

print("solveset(Eq(atan2(x, 1), pi), x, S.Reals) =", sol)
print("residual at returned value 0 =", (eq.lhs - eq.rhs).subs(x, 0))
print("atan2(0, 1) =", atan2(0, 1))
\end{lstlisting}

\subsection{Observed Behavior}

\medskip

\begin{lstlisting}[style=bugpython]
SymPy version: 1.14.0
solveset(Eq(atan2(x, 1), pi), x, S.Reals) = {0}
residual at returned value 0 = -pi
atan2(0, 1) = 0
\end{lstlisting}

SymPy reports \(x=0\) as a real solution. Direct substitution shows that this
point does not satisfy the equation: \(\operatorname{atan2}(0,1)=0\), so the
left-hand side minus the right-hand side is \(-\pi\).

\subsection{Expected Behavior}

The correct solution set is \(\emptyset\). For every real \(x\), the value
\(\operatorname{atan2}(x,1)\) is the principal argument of the point
\((1,x)\), which lies in the right half-plane. Its value is therefore strictly
between \(-\pi/2\) and \(\pi/2\), and it can never equal \(\pi\).

\subsection{Mathematical Explanation}

Because the second argument is positive,
\[
  \operatorname{atan2}(x,1)=\arctan(x)
\]
for real \(x\), using the principal real branch of \(\arctan\). This branch has
range
\[
  -\frac{\pi}{2}<\arctan(x)<\frac{\pi}{2}.
\]
The equation \(\operatorname{atan2}(x,1)=\pi\) is therefore equivalent to
\(\arctan(x)=\pi\), which has no real solution because \(\pi\) is outside the
range of \(\arctan\).

The returned set \(\{0\}\) is not a harmless branch representation or an
alternate form of the empty set. It is a false positive: substituting \(x=0\)
gives \(\operatorname{atan2}(0,1)=0\), not \(\pi\).

\subsection{Source-Level Diagnosis}

The source-level diagnosis identifies the relevant solver logic in
\nolinkurl{sympy/solvers/solveset.py}, especially the generic inverse-function
branch in \texttt{\_invert\_real} near lines 230--238 and the final checking
logic near lines 1395--1399. The traced call path is as follows. For a solve
over \(S.\mathrm{Reals}\), \texttt{solveset} first replaces the user symbol by a
real dummy in \nolinkurl{sympy/solvers/solveset.py}. Then
\(\operatorname{atan2}(\_R,1)\) evaluates to \(\arctan(\_R)\) in
\nolinkurl{sympy/functions/elementary/trigonometric.py}, because the second
argument is positive. The solver then handles
\(\arctan(\_R)-\pi=0\).

The problematic rule is the generic inversion step in \texttt{\_invert\_real}:

\begin{lstlisting}[style=bugpython]
imageset(Lambda(n, f.inverse()(n)), g_ys)
\end{lstlisting}

For \(f=\arctan\), this applies \(f^{-1}=\tan\) to the right-hand side and maps
\(\pi\) to \(\tan(\pi)=0\). This inversion is not reversible unless the
right-hand side is restricted to the actual real range of \(\arctan\), namely
\((-\pi/2,\pi/2)\). The diagnosis also notes that the later \texttt{\_check}
path only verifies definedness and finiteness through \texttt{domain\_check};
it does not verify that the equation residual is zero. As a result, the
spurious candidate \(0\) survives.

A plausible fix direction supported by the diagnosis is to intersect the
right-hand-side set with the range of the inverse trigonometric function before
using its \texttt{.inverse()} method. For \(\arctan\), values outside
\((-\pi/2,\pi/2)\) should be rejected before applying \(\tan\). A stronger
finite-solution residual check would also catch this class of false positives.

\subsection{Additional Failing Instances}

The checked cases cover the representative case and five additional instances
of the same pattern,
\(\operatorname{atan2}(a x,b)=\pm\pi\), with positive \(b\). In each case,
\(b>0\) places the point in the right half-plane, so the principal value is
again in \((-\pi/2,\pi/2)\). The equations with right-hand side \(\pi\) or
\(-\pi\) should all have empty real solution set.

\begin{center}
\small
\renewcommand{\arraystretch}{1.3}
\begin{tabular}{@{}>{\raggedright\arraybackslash}p{0.44\linewidth}>{\raggedright\arraybackslash}p{0.23\linewidth}>{\raggedright\arraybackslash}p{0.23\linewidth}@{}}
\toprule
\textbf{Input / Expression} & \textbf{SymPy behavior} & \textbf{Expected behavior} \\
\midrule
\(\operatorname{solveset}(\operatorname{atan2}(x,1)=\pi,\ x,\ \mathbb{R})\) & returns \(\{0\}\); residual at \(0\) is \(-\pi\) & \(\emptyset\) \\
\(\operatorname{solveset}(\operatorname{atan2}(x,2)=\pi,\ x,\ \mathbb{R})\) & returns \(\{0\}\); residual at \(0\) is \(-\pi\) & \(\emptyset\) \\
\(\operatorname{solveset}(\operatorname{atan2}(2x,1)=\pi,\ x,\ \mathbb{R})\) & returns \(\{0\}\); residual at \(0\) is \(-\pi\) & \(\emptyset\) \\
\(\operatorname{solveset}(\operatorname{atan2}(x/3,1)=\pi,\ x,\ \mathbb{R})\) & returns \(\{0\}\); residual at \(0\) is \(-\pi\) & \(\emptyset\) \\
\(\operatorname{solveset}(\operatorname{atan2}(-x,1)=\pi,\ x,\ \mathbb{R})\) & returns \(\{0\}\); residual at \(0\) is \(-\pi\) & \(\emptyset\) \\
\(\operatorname{solveset}(\operatorname{atan2}(x,1)=-\pi,\ x,\ \mathbb{R})\) & returns \(\{0\}\); residual at \(0\) is \(\pi\) & \(\emptyset\) \\
\bottomrule
\end{tabular}
\end{center}

\subsection{Related Issues Assessment}

The best-effort deduplication check found public issues in the same broad
family: SymPy issue \#6495 discusses \texttt{atan2} simplification and
inverse-trigonometric branch ambiguity, pull request \#21297 concerns improved
trigonometric solving in \texttt{solveset}, and issue \#12527 reports a
different \texttt{solveset} trigonometric domain-handling problem. None of the
reported matches identifies
\texttt{solveset(Eq(atan2(x, 1), pi), x, S.Reals)} returning \(\{0\}\), or the
specific false positive caused by reducing \(\operatorname{atan2}(x,1)\) to
\(\arctan(x)\) and applying \(\tan\) to an out-of-range right-hand side. The
deduplication verdict is therefore a new instance of a known failure mode: the
general branch/range issue is known, but this exact reproducible case remains
previously unreported and individually actionable as an issue and regression
test.

\subsection{Proposed SymPy Regression Test}

\medskip

\begin{lstlisting}[style=bugpython]
import pytest
from sympy import Eq, Rational, S, atan2, pi, solveset, symbols


x = symbols("x")


@pytest.mark.parametrize(
    "scale,denom,sign",
    [
        (S.One, S.One, S.One),
        (S.One, S(2), S.One),
        (S(2), S.One, S.One),
        (Rational(1, 3), S.One, S.One),
        (-S.One, S.One, S.One),
        (S.One, S.One, -S.One),
    ],
)
def test_solveset_atan2_rejects_values_outside_principal_range(scale, denom, sign):
    assert solveset(Eq(atan2(scale*x, denom), sign*pi), x, S.Reals) == S.EmptySet
\end{lstlisting}

\end{document}
